\chapter*{About this manual}
\addcontentsline{toc}{chapter}{About this manual}

This book describes TandemLLM version \version: the engine that serves Qwen3.8-27B on one NVIDIA DGX Spark, the conductor that sits in front of it, and the dashboard and proxy that watch it and other engines. It sits beside the Markdown pages in \code{docs/} and goes further: it shows how the parts fit together in 29 figures, and its reference part is generated from the code.

Part I says what TandemLLM is and follows one request through it. Part II installs and runs it, on a DGX Spark or as a dashboard on any machine. Part III explains how the engine reaches 49.89 tokens a second for one request, and how that number was measured. Part IV covers the conductor, the layer that plans each turn and watches it run. Part V walks through the dashboard's seven views. Part VI is the reference: about 700 rows of flags, variables, routes and metrics generated from the code, then troubleshooting and the roadmap. A glossary and an index close the book.

\section*{How numbers are labelled}

Every number in this book is one of three kinds, and the text says which.

\begin{itemize}
\item \measured: measured on one DGX Spark with the published NVFP4 weights, unless the text names another set. The caption or the sentence names the version or the date, and the tool that measured it when one did.
\item \arith: computed from sizes, such as bytes read divided by memory bandwidth. Nothing ran.
\item estimate: a guess with its basis stated. The book uses few of them and marks each one.
\end{itemize}

The speed numbers of the StairCut paper are the reference for single-request speed. They were measured on the published weights with the persistent suffix store off, on the benchmark row described in Chapter~\ref{ch:measurement}. Later releases are held to those numbers by a same-hour comparison before they ship.

Agreement with greedy decoding is a measured property with one documented exception, a near tie between the two best tokens. Chapter~\ref{ch:exactness} states it precisely. Where this book says that resuming from a cache is \enquote{bit for bit}, a test compared the bytes.

\section*{Conventions}

Commands, flags, variables, files and HTTP routes are set in \code{monospace}. Paths under \code{~/.qwen38-spark-engine/} and \code{~/.tandemllm/} are the engine's and the proxy's state directories on the machine that runs them. Runtime names still carry the project's old prefix (\code{QWEN38_}, \code{QSE_}); a later release renames them together.

The engine runs on Linux. The dashboard and the proxy run on Linux and macOS.

\section*{Building this book}

The source is in \code{manual/} of the repository. \code{make} in that directory regenerates the reference tables from the code (\code{tools/gen_reference.py}) and runs \code{latexmk -lualatex}. The repository's \code{manual.yml} workflow builds the same PDF on every change to \code{manual/} and attaches it to each tagged release.

\section*{License}\index{license}\index{AGPL-3.0-only}\index{CC BY 4.0}

The text of this manual is licensed under CC BY 4.0, like the documentation in \code{docs/}. The code it describes is licensed separately: AGPL-3.0-only, or a commercial license from the copyright holder. The model weights are not part of the repository; the base model and the base drafter are Apache-2.0, and so are the published NVFP4 overlays and fine-tuned drafters derived from them.

\vfill
{\small\color{slate} TandemLLM \version. \textcopyright{} 2026 the TandemLLM authors.}
